\documentclass[11pt,a4paper]{article}
\usepackage[T1]{fontenc}
\usepackage[utf8]{inputenc}
\usepackage{mathptmx}
\usepackage{amsmath,amssymb}
\usepackage{graphicx}
\usepackage{listings}
\usepackage{xcolor}
\usepackage{orcidlink}

% ---- Rd.sty FIRST, so our layout below always wins -----------------
\usepackage{Rd}
\providecommand{\proglang}[1]{\textsf{#1}}
\providecommand{\pkg}[1]{\textbf{#1}}
\providecommand{\CRANpkg}[1]{\href{https://cran.r-project.org/package=#1}{\pkg{#1}}}

% ---- page geometry: applied AFTER Rd.sty so it is never overridden -
\usepackage[margin=2.6cm]{geometry}
\usepackage[expansion=false]{microtype}
\emergencystretch=2em
\sloppy
\hypersetup{hidelinks}
\usepackage{titlesec}
\usepackage{fancyhdr}

% ---- look & feel --------------------------------------------------
\titleformat{\section}{\Large\bfseries}{\thesection}{1em}{}
\titleformat{\subsection}{\large\bfseries}{\thesubsection}{1em}{}
\setlength{\parskip}{0.4em}
\setlength{\parindent}{0em}

\pagestyle{fancy}
\fancyhf{}
\fancyhead[L]{\small\itshape The BvChen Package}
\fancyhead[R]{\small\itshape Bivariate Chen Distribution in R}
\fancyfoot[C]{\thepage}
\renewcommand{\headrulewidth}{0.4pt}

% ---- code listing style -------------------------------------------
\lstdefinestyle{bvchex}{
  basicstyle=\ttfamily\footnotesize,
  breaklines=true,
  columns=fullflexible,
  frame=single,
  framerule=0.3pt,
  rulecolor=\color{gray!60},
  backgroundcolor=\color{gray!7},
  xleftmargin=1em
}
\lstset{style=bvchex}
\begin{document}

% ======================= TITLE PAGE ===============================
\begin{titlepage}
\centering
\vspace*{1.5cm}
{\Huge \bfseries The Bivariate Chen Distribution\\[0.6em]
The \textsf{BvChen} R Package\\[1.2em]}
{\Large A toolkit for density, distribution, survival, dependence
measures, EM-based estimation and goodness-of-fit analysis\\[2em]}
{\large
\textbf{Mukul Bijalwan}\textsuperscript{1} \quad and \quad
\textbf{Puneet Kumar Gupta}\textsuperscript{2}\\[1em]}
{\normalsize
\textsuperscript{1}\ \texttt{mukulbijalwan555@gmail.com}\quad
\href{https://orcid.org/0009-0001-3040-6912}{%
\orcidlink{0009-0001-3040-6912}}\\
\textsuperscript{2}\ \texttt{puneetstat999@gmail.com}\\[2em]}
{\normalsize \today}
\vspace{1.5em}

\begin{minipage}{0.88\textwidth}
\begin{abstract}
\noindent
This manual documents the \textsf{BvChen} package, an \proglang{R}
implementation of the bivariate extension of the bathtub-shaped Chen
distribution proposed by Gupta, Pundir, Sharma and Mesfioui (2022,
\emph{Life Cycle Reliability and Safety Engineering} 11:247--259),
constructed from three independent latent Chen (2000) components
through competing risks. The package provides the joint density, distribution and
survival functions (including the singular component on the diagonal),
random variate generation, marginal and conditional distributions,
a Marshall--Olkin type survival copula with Kendall's $\tau$, joint
moments, maximum-likelihood estimation via the EM algorithm,
Kolmogorov--Smirnov goodness-of-fit tests and Q--Q diagnostics.
Complete univariate Chen building blocks are included. Every documented
function is illustrated with executable examples whose live output
appears directly beneath the code in this document.
\end{abstract}
\end{minipage}

\vspace{1em}
{\small \textbf{Keywords:} Chen distribution; bivariate lifetime
model; Marshall--Olkin structure; competing risks; EM algorithm;
survival copula; \proglang{R}.}

\vspace{1em}
{\footnotesize Generated automatically from the roxygen2 documentation
of version 0.1.0; all examples were executed while compiling this
manual.}
\vspace*{\fill}
\end{titlepage}

% ======================= TOC =======================================
\tableofcontents
\newpage

% ======================= THEORY ====================================
\section{The univariate Chen distribution}
The two-parameter Chen (2000) distribution has survival function
\begin{equation}
S(x)=\exp\{-\alpha(e^{x^{\beta}}-1)\}, \qquad x>0,
\; \alpha,\beta>0,
\label{eq:chen}
\end{equation}
hazard rate $h(x)=\alpha\beta x^{\beta-1}e^{x^{\beta}}$
(bathtub-shaped or increasing) and quantile function
$x_p=[\ln\{1-\ln(1-p)/\alpha\}]^{1/\beta}$.

\section{The bivariate construction}
Let $X_0,X_1,X_2$ be independent with $X_i \sim Ch(\alpha_i,\beta)$
and set $Z_1=\min(X_0,X_1)$, $Z_2=\min(X_0,X_2)$. Then
$Z_1 \sim Ch(\alpha_1+\alpha_3,\beta)$ and
$Z_2 \sim Ch(\alpha_2+\alpha_3,\beta)$ marginally. The joint law
is a mixture: an absolutely continuous part on $\{z_1<z_2\}$ and
$\{z_1>z_2\}$ carrying total mass
$(\alpha_1+\alpha_2)/(\alpha_1+\alpha_2+\alpha_3)$, and a
singular component on the diagonal $z_1=z_2$ of mass
$\alpha_3/(\alpha_1+\alpha_2+\alpha_3)$.

\subsection{Joint survival}
With $t(z)=e^{z^{\beta}}-1$,
\begin{equation}
S(z_1,z_2)=\exp\{-\alpha_1 t(z_1)-\alpha_2 t(z_2)
-\alpha_3\max(t(z_1),t(z_2))\}.
\label{eq:surv}
\end{equation}

\subsection{Survival copula}
In terms of marginal survival levels $u=S_1(z_1)$, $v=S_2(z_2)$,
\begin{equation}
\hat C(u,v)=\min\Bigl\{u^{\frac{\alpha_1}{\alpha_1+
\alpha_3}}v,\; uv^{\frac{\alpha_2}{\alpha_2+\alpha_3}}
\Bigr\},
\label{eq:cop}
\end{equation}
a Marshall--Olkin type copula whose dependence increases with the
relative size of the common-shock parameter $\alpha_3$ and is free
of $\beta$.

\subsection{Estimation by the EM algorithm}
On the exponential scale $t_i=e^{z_i^{\beta}}-1$ the latent model
reduces to a Marshall--Olkin competing-risks scheme. The package
alternates a numerical M-step maximising the observed log-likelihood
in $(\alpha_1,\alpha_2,\alpha_3)$ at fixed $\beta$ with a
one-dimensional update of $\beta$; observed ties $z_{1i}=z_{2i}$
identify the common shock exactly.

\newpage
% ==================== FUNCTION DOCUMENTATION =======================
\addcontentsline{toc}{section}{Function documentation}
\begin{center}{\Large\bfseries Function documentation}\end{center}
\vspace{0.5em}

\HeaderA{dbvch}{The Bivariate Chen Distribution}{dbvch}
\aliasA{pbvch}{dbvch}{pbvch}
\aliasA{rbvch}{dbvch}{rbvch}
\aliasA{sbvch}{dbvch}{sbvch}
%
\begin{Description}
Density (\code{dbvch}), joint CDF (\code{pbvch}),
joint survival function (\code{sbvch}) and random generation
(\code{rbvch}) for the bivariate Chen distribution built from
three independent latent Chen components: \eqn{Z_1=\min(X_0,X_1)}{},
\eqn{Z_2=\min(X_0,X_2)}{} with \eqn{X_i \sim Ch(\alpha_i,\beta)}{}.
\end{Description}
%
\begin{Usage}
\begin{verbatim}
dbvch(
  x,
  y,
  alpha1,
  alpha2,
  alpha3,
  beta,
  log = FALSE,
  component = c("full", "ac", "sing")
)

sbvch(x, y, alpha1, alpha2, alpha3, beta)

pbvch(x, y, alpha1, alpha2, alpha3, beta)

rbvch(n, alpha1, alpha2, alpha3, beta)
\end{verbatim}
\end{Usage}
%
\begin{Arguments}
\begin{ldescription}
\item[\code{x}, \code{y}] vectors of non-negative quantiles; recycled to common length.

\item[\code{alpha1}, \code{alpha2}, \code{alpha3}] positive shape parameters of the three
latent Chen variables.

\item[\code{beta}] positive common power parameter.

\item[\code{log}] logical; if TRUE the log-density of the continuous part
is returned (\code{-Inf} on the diagonal).

\item[\code{component}] character selector for \code{dbvch}: \code{"full"}
(default) evaluates the singular density on the diagonal and the
continuous density elsewhere; \code{"ac"} evaluates only the
absolutely continuous part (zero on the diagonal); \code{"sing"}
evaluates only the diagonal singular density.

\item[\code{n}] number of observations for \code{rbvch}.
\end{ldescription}
\end{Arguments}
%
\begin{Details}
The joint law is a mixture of an absolutely continuous
part living on \eqn{\{z_1<z_2\}}{} and \eqn{\{z_1>z_2\}}{}, carrying
total mass \eqn{(\alpha_1+\alpha_2)/(\alpha_1+\alpha_2+\alpha_3)}{},
and a singular part concentrated on the diagonal \eqn{z_1=z_2}{}
with mass \eqn{\alpha_3/(\alpha_1+\alpha_2+\alpha_3)}{}.  On
\eqn{z_1<z_2}{} the density factorises as
\eqn{f_{Ch(\alpha_1,\beta)}(z_1)\,
       f_{Ch(\alpha_2+\alpha_3,\beta)}(z_2)}{}.
\end{Details}
%
\begin{Value}
\begin{ldescription}
\item[\code{dbvch}] numeric vector of density values.
\item[\code{pbvch}] numeric vector, \eqn{P(Z_1\le x, Z_2\le y)}{}.
\item[\code{sbvch}] numeric vector, \eqn{P(Z_1>x, Z_2>y)}{}.
\item[\code{rbvch}] an \eqn{n\times 2}{} numeric matrix with columns
\code{z1}, \code{z2}, class \code{"bvch"}.
\end{ldescription}
\end{Value}
%
\begin{References}
Chen, Z. (2000). A new two-parameter lifetime distribution with
bathtub shape or increasing failure rate function.
\emph{Statistics \& Probability Letters} 49, 155-161.
Meintanis, S. G. (2007). Test of fit for Marshall-Olkin
distributions with applications. \emph{J. Statist. Comput. Simul.}
77, 171-179.
\end{References}
%


\vspace{0.8em}
\noindent{\bfseries Examples} (live output):
\begin{lstlisting}
> 
> 
> # density pieces
> dbvch(0.5, 1.5, 1, 1, 1, 1.5)      # region x < y
[1] 0.0005931645
> dbvch(1.5, 0.5, 1, 1, 1, 1.5)      # region x > y
[1] 0.0005931645
> dbvch(1.0, 1.0, 1, 1, 1, 1.5)      # singular diagonal density g(1)
[1] 0.02353232
> dbvch(1.0, 1.0, 1, 1, 1, 1.5, component = "ac")   # 0 on diagonal
[1] 0
> 
> # CDF from survival via inclusion-exclusion
> s  <- sbvch(1, 2, 1, 1, 1, 1.5)
> s1 <- schen(1, 2, 1.5); s2 <- schen(2, 2, 1.5)
> all.equal(pbvch(1, 2, 1, 1, 1, 1.5), 1 - s1 - s2 + s)
[1] TRUE
> 
> set.seed(1)
> sim <- rbvch(500, 1, 1, 1, 1.5)
> head(sim)
             z1        z2
[1,] 0.10617519 0.6814787
[2,] 0.08499553 0.1416199
[3,] 0.26448766 0.2644877
[4,] 0.25773862 0.2577386
[5,] 0.47240406 0.3201278
[6,] 1.22531700 0.2833057
> mean(sim[, "z1"] == sim[, "z2"])   # about a3/(a1+a2+a3) = 1/3
[1] 0.338
> 
> 
> 
\end{lstlisting}


\newpage

\HeaderA{dchen}{The Univariate Chen Distribution}{dchen}
\aliasA{hchen}{dchen}{hchen}
\aliasA{pchen}{dchen}{pchen}
\aliasA{qchen}{dchen}{qchen}
\aliasA{rchen}{dchen}{rchen}
\aliasA{schen}{dchen}{schen}
%
\begin{Description}
Density, distribution function, quantile function,
random generation, hazard function and maximum-likelihood fitting
for the Chen (2000) distribution with shape parameter \code{alpha}
and power parameter \code{beta}.
\end{Description}
%
\begin{Usage}
\begin{verbatim}
dchen(x, alpha, beta, log = FALSE)

pchen(q, alpha, beta)

schen(q, alpha, beta)

qchen(p, alpha, beta)

rchen(n, alpha, beta)

hchen(x, alpha, beta)
\end{verbatim}
\end{Usage}
%
\begin{Arguments}
\begin{ldescription}
\item[\code{x}, \code{q}] vector of quantiles (non-negative).

\item[\code{alpha}] shape parameter (\eqn{\alpha > 0}{}).

\item[\code{beta}] power parameter (\eqn{\beta > 0}{}).

\item[\code{log}] logical; if TRUE, the log-density is returned.

\item[\code{p}] vector of probabilities.

\item[\code{n}] number of observations.
\end{ldescription}
\end{Arguments}
%
\begin{Value}
\code{dchen} gives the density, \code{pchen} the CDF,
\code{schen} the survival function, \code{qchen} the quantile
function, \code{rchen} generates random deviates, \code{hchen}
gives the hazard rate and \code{fitchen} returns a fitted list
with elements \code{alpha}, \code{beta} and \code{loglik}.
\end{Value}
%
\begin{References}
Chen, Z. (2000). A new two-parameter lifetime distribution
with bathtub shape or increasing failure rate function.
\emph{Statistics \& Probability Letters} 49, 155--161.
\end{References}
%


\vspace{0.8em}
\noindent{\bfseries Examples} (live output):
\begin{lstlisting}
> 
> 
> dchen(1, alpha = 2, beta = 1.5)
[1] 0.2623826
> pchen(c(0.5, 1, 2), 2, 1.5)
[1] 0.5718313 0.9678249 1.0000000
> schen(1, 2, 1.5)                       # == 1 - pchen(1, 2, 1.5)
[1] 0.03217506
> qchen(c(0.25, 0.5, 0.75), 2, 1.5)
[1] 0.2623708 0.4457105 0.6521005
> hchen(seq(0.1, 2, length.out = 5), 2, 1.5)
[1]  0.9791627  3.5181490  9.0154268 24.3576403 71.7805109
> set.seed(1); x <- rchen(500, 2, 1.5)
> fit <- fitchen(x)                      # MLE close to (2, 1.5)
> fit$alpha; fit$beta
   alpha 
2.219127 
    beta 
1.611988 
> 
> 
> 
\end{lstlisting}


\newpage

\HeaderA{fitchen}{Univariate Maximum-Likelihood Fit of the Chen Distribution}{fitchen}
%
\begin{Description}
Fits the two-parameter Chen distribution by numerical
maximisation of the log-likelihood. Used internally to obtain
starting values for \code{\LinkA{fitbvch}{fitbvch}}.
\end{Description}
%
\begin{Usage}
\begin{verbatim}
fitchen(x, start = NULL, tol = 1e-08, maxit = 100)
\end{verbatim}
\end{Usage}
%
\begin{Arguments}
\begin{ldescription}
\item[\code{x}] non-negative numeric vector of observations.

\item[\code{start}] optional named list with starting values \code{alpha},
\code{beta}; sensible defaults are chosen from the data.

\item[\code{tol}] convergence tolerance used by \code{fitchen}.

\item[\code{maxit}] maximum number of Newton iterations in \code{fitchen}.
\end{ldescription}
\end{Arguments}
%
\begin{Value}
A list with components \code{alpha}, \code{beta},
\code{loglik}, \code{converged} and \code{se} (approximate
standard errors from the numerical Hessian).
\end{Value}
%


\vspace{0.8em}
\noindent{\bfseries Examples} (live output):
\begin{lstlisting}
> 
> 
> set.seed(42)
> x <- rchen(300, alpha = 1.5, beta = 2)
> fit <- fitchen(x)
> fit$alpha; fit$beta; fit$loglik
   alpha 
1.315234 
    beta 
2.028452 
[1] -24.91924
> 
> 
> 
\end{lstlisting}


\newpage

\HeaderA{Marginals}{Marginal Distributions of the Bivariate Chen Law}{Marginals}
\aliasA{dmarg1}{Marginals}{dmarg1}
\aliasA{dmarg2}{Marginals}{dmarg2}
\aliasA{pmarg1}{Marginals}{pmarg1}
\aliasA{pmarg2}{Marginals}{pmarg2}
%
\begin{Description}
Density and CDF of the marginals
\eqn{Z_1 \sim Ch(\alpha_1+\alpha_3,\beta)}{} and
\eqn{Z_2 \sim Ch(\alpha_2+\alpha_3,\beta)}{}.
\end{Description}
%
\begin{Usage}
\begin{verbatim}
dmarg1(x, alpha1, alpha2, alpha3, beta)

dmarg2(y, alpha1, alpha2, alpha3, beta)

pmarg1(q, alpha1, alpha2, alpha3, beta)

pmarg2(q, alpha1, alpha2, alpha3, beta)
\end{verbatim}
\end{Usage}
%
\begin{Arguments}
\begin{ldescription}
\item[\code{x}, \code{q}] non-negative quantiles.

\item[\code{alpha1}, \code{alpha2}, \code{alpha3}] positive shape parameters of the three
latent Chen variables.

\item[\code{beta}] positive common power parameter.

\item[\code{y}] quantile at which the second marginal is evaluated.
\end{ldescription}
\end{Arguments}
%
\begin{Value}
Numeric vectors of density (\code{d*}) or probability
(\code{p*}) values.
\end{Value}
%


\vspace{0.8em}
\noindent{\bfseries Examples} (live output):
\begin{lstlisting}
> 
> 
> dmarg1(1, 1, 1, 1, 1.5)
[1] 0.2623826
> dchen(1, 2, 1.5)                   # identical by construction
[1] 0.2623826
> pmarg2(c(0.5, 1, 2), 1, 1, 1, 1.5)
[1] 0.5718313 0.9678249 1.0000000
> pchen(c(0.5, 1, 2), 2, 1.5)
[1] 0.5718313 0.9678249 1.0000000
> 
> 
> 
\end{lstlisting}


\newpage

\HeaderA{Conditional}{Conditional Distributions of the Bivariate Chen Law}{Conditional}
\aliasA{dcond1}{Conditional}{dcond1}
\aliasA{dcond2}{Conditional}{dcond2}
\aliasA{hcond1}{Conditional}{hcond1}
\aliasA{hcond2}{Conditional}{hcond2}
%
\begin{Description}
Conditional density (\code{dcond}) and conditional
hazard (\code{hcond}) of \eqn{Z_1}{} given \eqn{Z_2=y}{} (suffix 1)
and of \eqn{Z_2}{} given \eqn{Z_1=x}{} (suffix 2).
\end{Description}
%
\begin{Usage}
\begin{verbatim}
dcond1(z1, y, alpha1, alpha2, alpha3, beta)

dcond2(z2, x, alpha1, alpha2, alpha3, beta)

hcond1(z1, y, alpha1, alpha2, alpha3, beta)

hcond2(z2, x, alpha1, alpha2, alpha3, beta)
\end{verbatim}
\end{Usage}
%
\begin{Arguments}
\begin{ldescription}
\item[\code{z1}, \code{z2}] values at which the conditional density/hazard of the
response variable is evaluated (any non-negative value).

\item[\code{alpha1}, \code{alpha2}, \code{alpha3}] positive shape parameters of the three
latent Chen variables.

\item[\code{beta}] positive common power parameter.

\item[\code{x}, \code{y}] conditioning values (\code{x} conditions on
\eqn{Z_1=x}{}, \code{y} on \eqn{Z_2=y}{}); must be positive.
\end{ldescription}
\end{Arguments}
%
\begin{Details}
Given \eqn{Z_2=y}{}, the conditional law of \eqn{Z_1}{} has an
atom at \eqn{z_1=y}{} of size \eqn{\pi(y)=g(y)/f_{m2}(y)}{}, where
\eqn{g}{} is the diagonal singular density.  Its continuous part is
piecewise:
\begin{itemize}

\item{} \eqn{z_1<y}{}: \eqn{f_{Ch(\alpha_1,\beta)}(z_1)}{};
\item{} \eqn{z_1>y}{}: \eqn{f_{Ch(\alpha_2,\beta)}(y)\,
          f_{Ch(\alpha_1+\alpha_3,\beta)}(z_1)/f_{m2}(y)}{},

\end{itemize}

with \eqn{f_{m2}}{} the marginal density of \eqn{Z_2}{}.

The conditional hazard of the absolutely continuous part is
\eqn{h(z\mid y)=f_{Z_1|Z_2}(z\mid y)/P(Z_1>z\mid Z_2=y)}{}; it is
\code{NA} exactly at \eqn{z=y}{}, where only the point mass lives.
\end{Details}
%
\begin{Value}
Numeric vector with attribute \code{"atom"} giving the
conditional probability mass at \eqn{z_i = y}{}; the returned
density itself refers to the strictly-off-diagonal part.
\code{hcond*} return the conditional hazard rate of the
absolutely continuous part (\code{NA} exactly at \eqn{z=y}{}).
\end{Value}
%


\vspace{0.8em}
\noindent{\bfseries Examples} (live output):
\begin{lstlisting}
> 
> 
> # conditional density of Z1 given Z2 = 1, below and above y = 1
> round(dcond1(seq(0.4, 1.8, by = 0.2), 1, 1, 1, 1, 1.5), 4)
[1] 0.9162 1.0235 0.9648 0.7314 0.1471 0.0107 0.0002 0.0000
attr(,"atom")
[1] 0.08968704
> attr(dcond1(1, 1, 1, 1, 1, 1.5), "atom")   # point mass at z1 = y
[1] 0.08968704
> 
> round(hcond1(seq(0.4, 1.8, by = 0.2), 1, 1, 1, 1, 1.5), 4)
[1]  1.2218  1.8493  2.7441      NA 10.9045 16.5804 25.5936 40.1383
> 
> 
\end{lstlisting}


\newpage

\HeaderA{tau\_bvch}{Dependence Measures of the Bivariate Chen Law}{tau.Rul.bvch}
\aliasA{concordance\_bvch}{tau\_bvch}{concordance.Rul.bvch}
\aliasA{Dependence}{tau\_bvch}{Dependence}
%
\begin{Description}
Kendall's tau (Monte-Carlo estimate via
\eqn{\tau = 4\,E[F(Z_1,Z_2)] - 1}{}) and a concordance comparison of
two parameter triples.  Dependence increases with the relative
size of the common-shock parameter \eqn{\alpha_3}{} and does not
depend on \eqn{\beta}{}.
\end{Description}
%
\begin{Usage}
\begin{verbatim}
tau_bvch(alpha1, alpha2, alpha3, beta, n = 50000, seed = NULL)

concordance_bvch(
  alpha1a,
  alpha2a,
  alpha3a,
  alpha1b,
  alpha2b,
  alpha3b,
  beta,
  n = 50000,
  seed = NULL
)
\end{verbatim}
\end{Usage}
%
\begin{Arguments}
\begin{ldescription}
\item[\code{alpha1}, \code{alpha2}, \code{alpha3}] positive shape parameters.

\item[\code{beta}] positive power parameter.

\item[\code{n}] number of Monte-Carlo draws for \code{tau\_bvch}.

\item[\code{seed}] optional seed for reproducibility.

\item[\code{alpha1a}, \code{alpha2a}, \code{alpha3a}] first parameter triple of
\code{concordance\_bvch}.

\item[\code{alpha1b}, \code{alpha2b}, \code{alpha3b}] second parameter triple of
\code{concordance\_bvch}.
\end{ldescription}
\end{Arguments}
%
\begin{Value}
\code{tau\_bvch}: numeric estimate of Kendall's tau.
\code{concordance\_bvch}: TRUE if the first triple is
concordance-smaller than the second.
\end{Value}


\newpage

\HeaderA{Moments}{Moments of the (Bivariate) Chen Distribution}{Moments}
\aliasA{ebvch}{Moments}{ebvch}
\aliasA{mbvch}{Moments}{mbvch}
\aliasA{mchen}{Moments}{mchen}
\aliasA{min\_bvch}{Moments}{min.Rul.bvch}
\aliasA{vbvch}{Moments}{vbvch}
%
\begin{Description}
Moments of the (Bivariate) Chen Distribution
\end{Description}
%
\begin{Usage}
\begin{verbatim}
mchen(k, alpha, beta)

mbvch(r, s, alpha1, alpha2, alpha3, beta, tol = 1e-09)

ebvch(alpha1, alpha2, alpha3, beta)

vbvch(alpha1, alpha2, alpha3, beta)

min_bvch(alpha1, alpha2, alpha3, beta)
\end{verbatim}
\end{Usage}
%
\begin{Arguments}
\begin{ldescription}
\item[\code{k}] moment order for the univariate integral.

\item[\code{alpha}] shape parameter (\eqn{\alpha > 0}{}) of the univariate
Chen distribution.

\item[\code{beta}] positive common power parameter.

\item[\code{r}, \code{s}] non-negative orders of the moments \eqn{E[Z_1^r]}{},
\eqn{E[Z_2^s]}{} and the cross moment \eqn{E[Z_1^r Z_2^s]}{}.

\item[\code{alpha1}, \code{alpha2}, \code{alpha3}] positive shape parameters of the three
latent Chen variables.

\item[\code{tol}] stopping tolerance for the numerical integration.
\end{ldescription}
\end{Arguments}
%
\begin{Value}
\begin{ldescription}
\item[\code{mchen}] numeric, \eqn{E[Z^k]}{} for \eqn{Z~Ch(alpha,beta)}{}.
\item[\code{mbvch}] numeric, cross moment \eqn{E[Z_1^r Z_2^s]}{}.
\item[\code{ebvch}] named vector of means \code{c(EZ1, EZ2)}.
\item[\code{vbvch}] 2x2 variance-covariance matrix.
\item[\code{min\_bvch}] list describing the minimum
\eqn{M=\min(Z_1,Z_2)\sim Ch(\alpha_1+\alpha_2+\alpha_3,\beta)}{}
with elements \code{alpha}, \code{beta}, and functions
\code{d}, \code{p}, \code{q}.
\end{ldescription}
\end{Value}
%


\vspace{0.8em}
\noindent{\bfseries Examples} (live output):
\begin{lstlisting}
> 
> 
> mchen(1, 2, 1.5)                       # E[Z], Ch(2, 1.5)
[1] 0.4707204
> mean(rchen(200000, 2, 1.5))            # agrees up to MC error
[1] 0.4701932
> 
> mbvch(1, 1, 1, 1, 1, 1.5)
[1] 0.2498289
> ebvch(1, 1, 1, 1.5)
      EZ1       EZ2 
0.4707204 0.4707204 
> vbvch(1, 1, 1, 1.5)
           Z1         Z2
Z1 0.06964704 0.02825116
Z2 0.02825116 0.06964704
> 
> mn <- min_bvch(1, 1, 1, 1.5)
> mn$p(1)                                # P(min <= 1)
[1] 0.9942286
> pchen(1, 3, 1.5)                       # identical
[1] 0.9942286
> 
> 
\end{lstlisting}


\newpage

\HeaderA{Likelihood}{Log-Likelihood and Profile Likelihood of the Bivariate Chen Model}{Likelihood}
\aliasA{loglik\_bvch}{Likelihood}{loglik.Rul.bvch}
\aliasA{profile\_beta}{Likelihood}{profile.Rul.beta}
%
\begin{Description}
Observed-data log-likelihood and profile likelihood for
the common power parameter \eqn{\beta}{}.
\end{Description}
%
\begin{Usage}
\begin{verbatim}
profile_beta(
  z1,
  z2,
  betas = NULL,
  start = c(alpha1 = 1, alpha2 = 1, alpha3 = 1)
)

loglik_bvch(z1, z2, alpha1, alpha2, alpha3, beta)
\end{verbatim}
\end{Usage}
%
\begin{Arguments}
\begin{ldescription}
\item[\code{z1}, \code{z2}] numeric vectors of paired observations
(\code{z1[i]}, \code{z2[i]}).

\item[\code{betas}] candidate beta values for the profile grid; a sensible
default grid over (0, 20] is used if omitted.

\item[\code{start}] named numeric vector with starting values for the
alpha search inside \code{profile\_beta}.

\item[\code{alpha1}, \code{alpha2}, \code{alpha3}] positive shape parameters of the three
latent Chen variables.

\item[\code{beta}] positive common power parameter.
\end{ldescription}
\end{Arguments}
%
\begin{Value}
\code{loglik\_bvch} returns the scalar log-likelihood;
\code{profile\_beta} returns a list with the evaluated grid
(\code{beta}, \code{prof\_loglik}) and the optimiser
\code{beta\_hat}.
\end{Value}
%


\vspace{0.8em}
\noindent{\bfseries Examples} (live output):
\begin{lstlisting}
> 
> 
> set.seed(7)
> sim <- rbvch(300, 1, 1, 1, 1.5)
> loglik_bvch(sim[, "z1"], sim[, "z2"], 1, 1, 1, 1.5)
[1] -205.452
> 
> pb <- profile_beta(sim[, "z1"], sim[, "z2"])
> pb$beta_hat                        # close to 1.5
[1] 1.523276
> head(pb$grid)
        beta prof_loglik
1 0.05000000  -1508.2561
2 0.06417844  -1384.5584
3 0.08237745  -1261.8931
4 0.10573713  -1140.5872
5 0.13572088  -1021.0718
6 0.17420709   -903.9121
> 
> 
\end{lstlisting}


\newpage

\HeaderA{fitbvch}{EM Estimation of the Bivariate Chen Distribution}{fitbvch}
\aliasA{coef.bvchfit}{fitbvch}{coef.bvchfit}
\aliasA{print.bvchfit}{fitbvch}{print.bvchfit}
%
\begin{Description}
Fits the four-parameter bivariate Chen distribution
\eqn{(\alpha_1,\alpha_2,\alpha_3,\beta)}{} to paired data with the
EM algorithm.  On the exponential scale \eqn{t=e^{z^\beta}-1}{} the
model is a Marshall-Olkin competing-risks scheme, so each
iteration re-estimates the \eqn{\alpha}{}'s in closed form
(expected failures over exposure) and updates \eqn{\beta}{} by
one-dimensional numerical maximisation of the observed
log-likelihood.  An observed tie \eqn{z_{1i}=z_{2i}}{} identifies
the common-shock component exactly.
\end{Description}
%
\begin{Usage}
\begin{verbatim}
fitbvch(z1, z2, start = NULL, tol = 1e-08, maxit = 200L, verbose = FALSE)

## S3 method for class 'bvchfit'
coef(object, ...)

## S3 method for class 'bvchfit'
print(x, digits = max(3L, getOption("digits") - 3L), ...)
\end{verbatim}
\end{Usage}
%
\begin{Arguments}
\begin{ldescription}
\item[\code{z1}, \code{z2}] paired non-negative observations.

\item[\code{start}] optional named list with starting values
\code{alpha1}, \code{alpha2}, \code{alpha3}, \code{beta}; if
omitted, values derived from marginal univariate fits are used.

\item[\code{tol}] relative tolerance on the log-likelihood increase.

\item[\code{maxit}] maximum number of EM iterations.

\item[\code{verbose}] logical; print the iteration history.

\item[\code{object}] an object of class \code{"bvchfit"}.

\item[\code{...}] further arguments (unused).

\item[\code{x}] an object of class \code{"bvchfit"}.

\item[\code{digits}] number of digits passed to \code{\LinkA{round}{round}}.
\end{ldescription}
\end{Arguments}
%
\begin{Value}
An object of class \code{"bvchfit"}: list with elements
\begin{ldescription}
\item[\code{coefficients}] named vector \code{(alpha1, alpha2, alpha3,
    beta)}
\item[\code{loglik}] observed log-likelihood at the solution
\item[\code{iterations}] number of EM iterations performed
\item[\code{converged}] logical convergence flag
\item[\code{tau}] Kendall's tau implied by the fitted parameters
\item[\code{n}] sample size
\item[\code{n\_ties}] number of observed ties \eqn{z_{1i}=z_{2i}}{}
\item[\code{history}] iteration-wise log-likelihood values
\end{ldescription}
\end{Value}
%
\begin{References}
Meintanis, S. G. (2007). Test of fit for
Marshall-Olkin distributions with applications.  \emph{J. Statist.
Comput. Simul.} 77, 171-179.
\end{References}
%


\vspace{0.8em}
\noindent{\bfseries Examples} (live output):
\begin{lstlisting}
> 
> 
> set.seed(2024)
> sim <- rbvch(400, 0.8, 1.2, 0.5, 1.5)
> fit <- fitbvch(sim[, "z1"], sim[, "z2"])
> round(fit$coefficients, 3)          # near (0.8, 1.2, 0.5, 1.5)
alpha1 alpha2 alpha3   beta 
 0.776  1.368  0.645  1.572 
> fit$loglik
[1] -264.4934
> fit$tau                             # dependence level
[1] 0.2379021
> 
> # invariant to starting values
> fit2 <- fitbvch(sim[, "z1"], sim[, "z2"],
                start = list(alpha1 = 2, alpha2 = 2,
                             alpha3 = 2, beta = 1))
> all.equal(coef(fit), coef(fit2), tolerance = 1e-4)
[1] "Mean relative difference: 0.0001562438"
> 
> 
> 
\end{lstlisting}


\newpage

\HeaderA{gofbvch}{Goodness-of-Fit Tests and Q-Q Plots}{gofbvch}
\aliasA{qqbvch}{gofbvch}{qqbvch}
%
\begin{Description}
Kolmogorov-Smirnov goodness-of-fit tests for the two
marginals and for the minimum \eqn{\min(Z_1,Z_2)}{}, following the
testing strategy of Meintanis (2007), together with Q-Q plots for
\eqn{Z_1}{}, \eqn{Z_2}{} and the minimum.
\end{Description}
%
\begin{Usage}
\begin{verbatim}
gofbvch(z1, z2, alpha1, alpha2, alpha3, beta)

qqbvch(z1, z2, alpha1, alpha2, alpha3, beta, plot = TRUE)
\end{verbatim}
\end{Usage}
%
\begin{Arguments}
\begin{ldescription}
\item[\code{z1}, \code{z2}] paired non-negative observations.

\item[\code{alpha1}, \code{alpha2}, \code{alpha3}, \code{beta}] fitted parameter values.

\item[\code{plot}] logical; if TRUE (default) the Q-Q plots are drawn.
\end{ldescription}
\end{Arguments}
%
\begin{Value}
\code{gofbvch} returns a list with elements
\code{ks\_z1}, \code{ks\_z2}, \code{ks\_min}: each an object of
class \code{"htest"} produced by \code{\LinkA{ks.test}{ks.test}},
comparing the data against the corresponding theoretical CDF.

\code{qqbvch} invisibly returns the sorted data and theoretical
quantiles used in the plots, as a list with elements
\code{z1}, \code{z2}, \code{min}.
\end{Value}
%
\begin{References}
Meintanis, S. G. (2007). Test of fit for
Marshall-Olkin distributions with applications.
\emph{J. Statist. Comput. Simul.} 77, 171-179.
\end{References}


\newpage

\HeaderA{plots}{Plots for the Bivariate Chen Distribution}{plots}
\aliasA{persp.bvch}{plots}{persp.bvch}
%
\begin{Description}
Contour (\code{plot}) and perspective (\code{persp})
views of the absolutely continuous part of the joint density,
plus a scatter-plot method for simulated data of class
\code{"bvch"}.
\end{Description}
%
\begin{Usage}
\begin{verbatim}
## S3 method for class 'bvch'
persp(
  x,
  alpha1,
  alpha2,
  alpha3,
  beta,
  from = 0.05,
  to = 2.5,
  n.grid = 50,
  theta = -30,
  phi = 25,
  ...
)
\end{verbatim}
\end{Usage}
%
\begin{Arguments}
\begin{ldescription}
\item[\code{x}] for \code{plot.bvch}: an \code{n x 2} matrix of class
\code{"bvch"} (as produced by \code{\LinkA{rbvch}{rbvch}}) or an object
of class \code{"bvchfit"}; for \code{persp.bvch} the same types.

\item[\code{alpha1}, \code{alpha2}, \code{alpha3}, \code{beta}] distribution parameters (required
when \code{x} is a data matrix).

\item[\code{from}, \code{to}] range of the evaluation grid on both axes.

\item[\code{n.grid}] grid size per axis.

\item[\code{theta}, \code{phi}] viewing angles passed to \code{\LinkA{persp}{persp}}.

\item[\code{...}] further graphical parameters.
\end{ldescription}
\end{Arguments}
%
\begin{Value}
\code{plot.bvch} invisibly returns the evaluated density
grid; \code{persp.bvch} invisibly returns the transformation
matrix used.
\end{Value}
%


\vspace{0.8em}
\noindent{\bfseries Examples} (live output):
\begin{lstlisting}
> 
> 
> set.seed(11)
> sim <- rbvch(300, 1, 1, 1, 1.5)
> 
> # contour of the AC density with the simulated scatter
> plot(sim, 1, 1, 1, 1.5, to = 2.5)
> 
> # 3-D perspective of the absolutely continuous part
> persp(sim, 1, 1, 1, 1.5, to = 2)
> 
> 
\end{lstlisting}


\newpage
% ==================== REFERENCES ====================================
\newpage
\section*{References}
\addcontentsline{toc}{section}{References}
{\small
\noindent Ahmed EA (2014) Bayesian estimation based on progressive type-II censoring from two-parameter bathtub-shaped lifetime model: an Markov chain Monte Carlo approach. J Appl Stat 41(4):752--768\\[0.35em]
\noindent Bhatti FA, Hamedani GG, Najibi SM, Ahmad M (2021) On the extended Chen distribution: development, properties, characterizations and applications. Ann Data Sci 8(1):159--180\\[0.35em]
\noindent Block H, Basu AP (1974) A continuous bivariate exponential extension. J Am Stat Assoc 69:1031--1037\\[0.35em]
\noindent Chen Z (2000) A new two-parameter lifetime distribution with bathtub shape or increasing failure rate function. Stat Probab Lett 49(2):155--161\\[0.35em]
\noindent Dempster AP, Laird NM, Rubin DB (1977) Maximum likelihood from incomplete data via the EM algorithm. J R Stat Soc Ser B (Methodol) 39(1):1--38\\[0.35em]
\noindent Dey S, Sharma VK, Mesfioui M (2017) A new extension of Weibull distribution with application to lifetime data. Ann Data Sci. https://doi.org/10.1007/s40745-016-0094-8\\[0.35em]
\noindent Gupta PK, Pundir PS, Sharma VK, Mesfioui M (2022) Bivariate extension of bathtub-shaped distribution. Life Cycle Reliab Saf Eng 11:247--259. https://doi.org/10.1007/s41872-022-00193-4\\[0.35em]
\noindent Johnson NL, Kotz S, Kemp AW (1992) Univariate discrete distributions, 2nd edn. Wiley, New York\\[0.35em]
\noindent Johnson NL, Kotz S, Balakrishnan N (1995) Continuous univariate distributions, vol 2. Wiley, New York\\[0.35em]
\noindent Johnson NL, Kotz S, Balakrishnan N (1997) Discrete multivariate distributions, vol 165. Wiley, New York\\[0.35em]
\noindent Kayal T, Tripathi YM, Singh DP, Rastogi MK (2016) Estimation and prediction for Chen distribution with bathtub shape under progressive censoring. J Stat Comput Simul. https://doi.org/10.1080/00949655.2016.1209199\\[0.35em]
\noindent Kocherlakota S, Kocherlakota K (1992) Bivariate discrete distributions. Wiley, New York\\[0.35em]
\noindent Kotz S, Balakrishnan N, Johnson NL (2004) Continuous multivariate distributions, volume 1: models and applications. Wiley, New York\\[0.35em]
\noindent Kundu D, Dey AK (2009) Estimating the parameters of the Marshall--Olkin bivariate Weibull distribution by EM algorithm. Comput Stat Data Anal 53(4):956--965\\[0.35em]
\noindent Kundu D, Gupta RD (2009) Bivariate generalized exponential distribution. J Multivar Anal 100:581--593\\[0.35em]
\noindent Kundu D, Gupta RD (2010) Modified Sarhan--Balakrishnan singular bivariate distribution. J Stat Plan Inference 140(2):526--538\\[0.35em]
\noindent Marshall AW, Olkin I (1967) A multivariate exponential distribution. J Am Stat Assoc 62:30--44\\[0.35em]
\noindent Meintanis SG (2007) Test of fit for Marshall--Olkin distributions with applications. J Stat Plan Inference 137(12):3954--3963\\[0.35em]
\noindent Papageorgiou H (1997) Multivariate discrete distributions. In: Kotz S, Read CB, Banks DL (eds) Encyclopedia of statistical sciences, vol 1. Wiley, New York, pp 408--419\\[0.35em]
\noindent Pundir PS, Gupta PK (2018a) Reliability estimation in load-sharing system model with application to real data. Ann Data Sci 5(1):69--91\\[0.35em]
\noindent Pundir PS, Gupta PK (2018b) Stress-strength reliability of two-parameter bathtub-shaped lifetime model based on hybrid censored samples. J Stat Manag Syst 21(7):1229--1250\\[0.35em]
\noindent Rajarshi S, Rajarshi MB (1988) Bathtub distributions: a review. Commun Stat Theory Methods 17(8):2597--2621\\[0.35em]
\noindent Rastogi MK, Tripathi YM (2013) Estimation using hybrid censored data from a two parameter distribution with bathtub shape. Comput Stat Data Anal 67:268--281\\[0.35em]
\noindent Sarhan AM, Hamilton DC, Smith B (2012) Parameter estimations for a two parameter bathtub-shaped lifetime distribution. Appl Math Model 36(11):5380--5392\\[0.35em]
\noindent Tyagi A, Choudhary N, Singh B (2021) Reliability analysis of the dynamic system for the Chen model through sequential order statistics. Qual Reliab Eng Int 37(6):2514--2534\\[0.35em]
\noindent Wang L, Wu K, Tripathi YM, Lodhi C (2022) Reliability analysis of multicomponent stress--strength reliability from a bathtub-shaped distribution. J Appl Stat 49(1):122--142\\[0.35em]
\noindent Wu SJ (2008) Estimation of the two-parameter bathtub-shaped lifetime distribution with progressive censoring. J Appl Stat 35:1139--1150\\[0.35em]
\noindent Wu JW, Lu HL, Chen CH, Wu CH (2004) Statistical inference about the shape parameter of the new two-parameter bathtub-shaped lifetime distribution. Qual Reliab Eng Int 20:607--616
}

\end{document}

